\documentclass{NoteThesis/note-thesis}

\course{Advanced Polymer Chemistry}

\coursetitle{QXU6033 Advanced Polymer Chemistry}
\semester{2025B}
\teacher{Prof. Yamin Abdouni and Prof. Sulafudin Vukusic}

\begin{document}

\maketitle

\tableofcontents

\newpage

\section{Introduction to Module}

\subsection*{Assessment}

\begin{table}[h]
    \centering
    \begin{tabular}{cc}
        \hline
        Activity & Score\\
        \hline
        Professional Practice Evaluation(supervision) & 20\%\\
        Coursework & 30\%\\
        Final Exam & 50\%\\ 
        \hline
    \end{tabular}
\end{table}

\subsection{Study Goal}

\subsubsection*{Controlled Radical Polymerisation}
Mechanistic principles of controlled / living radical polymerisation (CRP)

Kinetic control, equilibrium processes, and radical deactivation strategies

Molecular weight control, dispersity (Đ), and end-group fidelity

\subsubsection*{Core CRP methodologies}

Nitroxide-Mediated Polymerisation (NMP)

Reversible Addition–Fragmentation Chain Transfer (RAFT)

Atom Transfer Radical Polymerisation (ATRP)

\subsubsection*{Advanced macromolecular architectures}

Self-Condensing Vinyl Polymerisation (SCVP)

Strathclyde synthesis and hyperbranched structures

Star polymers and multi-arm architectures

\subsubsection*{From synthesis to structure}

Block copolymer design and synthesis

Microphase separation and self-assembly

\begin{mynote_env}[frametitle=Introductory Discussions I]
    \textbf{Q:} A has a structure that enables it to polymerise through chain-growth polymerisation, while B is step-growth.\\ Scenario 1: Imagine you start polymerising A, and half-way through the polymerisation, you add C, also another monomer, which can polymerise through chain-growth polymerisation, and, that can polymerise together well with A (this is called “copolymerise”). Separately, imagine Scenario 2: Use B instead A, then add D, a monomer, which can polymerise through step-growth polymerisation, and, that can copolymerise well with B. At the end of both polymerisation processes, what type of polymer samples will you have in each Scenario?\\
    \textbf{A:} S1: AAAAA-(random A and C)-(if one is finish, pure another one).\\ S2:(block B)-some D-(block B) and cycle B.
\end{mynote_env}

\end{document}
